\documentclass{article}
\usepackage[margin=1in]{geometry}
\usepackage{amsmath}
\usepackage{amssymb}
\usepackage{bm}
\usepackage{amsthm}

\newtheorem{theorem}{Theorem}
\newtheorem{lemma}{Lemma}
\newtheorem{corollary}{Corollary}
\newtheorem{assumption}{Assumption}
\newtheorem{remark}{Remark}

\begin{document}

\title{Nonparametric Predictive Regression with Exogenous Shocks: A Sieve Empirical Likelihood Approach (In progress Version 2.0 AR1)}
\author{Siyan Dong}
\date{\today}
\maketitle

\section{Introduction and Data-Generating Process}

We consider the problem of testing return predictability in the presence of highly persistent predictors and exogenous covariates (such as weather shocks). The data-generating process (DGP) is specified by the following system of equations:

\begin{enumerate}
    \item \textbf{Predictive Regression:}
    \begin{equation}
        Y_t = m(W_{t-1}) + \beta X_{t-1} + U_t
    \end{equation}
    \begin{equation}
        X_t = \theta + \rho X_{t-1} + \varepsilon_t
    \end{equation}
    
    \item \textbf{Endogeneity:}
    \begin{equation}
        U_t = \phi \varepsilon_t + z_t \quad \text{with } E[z_t \mid \varepsilon_t] = 0
    \end{equation}
\end{enumerate}
where $Y_t$ represents the asset returns, $W_{t-1}$ is a stationary exogenous weather covariate, and $X_{t-1}$ is a persistent financial predictor that is independent of $W_{t-1}$. The error terms are structured to accommodate potential endogeneity via the parameter $\phi$.

Motivated by Perron (1989), our goal is to investigate return predictability while avoiding size distortions or false inferences caused by near unit roots in the predictor dynamics (i.e., $\rho \approx 1$). This is achieved by testing the null hypothesis of no predictability:
\begin{equation} 
\begin{aligned}
    H_0: \quad &\beta = 0 \\
    H_1: \quad &\beta \neq 0
\end{aligned}
\end{equation}

We propose a two-stage Orthogonalized Sieve-EL framework that achieves a standard $\chi^2_1$ limiting distribution under the null under high-level persistent-predictor score conditions.

\section{Empirical Motivation: Weather Shocks and Cocoa Futures}

We motivate our framework with the 2023--2024 agricultural commodity shocks. Let $Y_t$ denote excess returns on cocoa futures, $W_{t-1}$ an exogenous local weather variable such as precipitation or average temperature, and $X_{t-1}$ a persistent financial predictor, such as the lagged price. 

During 2023--2024, extreme weather conditions, including severe droughts in West Africa, led to supply deficits and a substantial increase in cocoa prices. This context illustrates the key features of our model:

\begin{enumerate}
    \item \textbf{Exogeneity of Weather Shocks:} Local weather variables ($W_{t-1}$) are determined by atmospheric systems, exogenous to financial variables ($X_{t-1}$) such as trading volume or lagged prices. Thus, $W_{t-1}$ and $X_{t-1}$ are independent, satisfying Assumption~A.
    
    \item \textbf{Nonlinear Weather Impact:} The effect of temperature and rainfall on crop yield is highly nonlinear, with extreme weather deviations causing disproportionate impacts on supply. A nonparametric formulation $m(W_{t-1})$ is necessary to capture this threshold behavior, which we approximate using a sieve basis.
    
    \item \textbf{Predictor Persistence:} Financial predictors ($X_{t-1}$) often exhibit near unit root behavior ($\rho \approx 1$). Standard OLS inference on $\beta$ with $m(W_{t-1})$ and persistent $X_{t-1}$ leads to size distortions.
    
    \item \textbf{Orthogonalized Sieve-EL Application:} Our framework first filters the nonlinear weather impact $m(W_{t-1})$ via sieve projection. The independence of $W_{t-1}$ and $X_{t-1}$ ensures consistency. Centering the predictor's weights against the sieve basis purges the nonparametric estimation error from the EL moment condition, allowing for standard $\chi^2_1$ inference on $\beta$.
\end{enumerate}

\section{The Orthogonalized Sieve-EL Framework (Version 2.0 AR1)}

We handle the unknown nonparametric curve $m(W_{t-1})$ by approximating it using a sieve basis of dimension $K$ (e.g., polynomials or splines). Let $\bm{P}_K(W_{t-1}) = (p_1(W_{t-1}), \dots, p_K(W_{t-1}))'$ be the $K \times 1$ vector of these basis functions.

\subsection{The Step-by-Step Algorithm}

The estimation and inference procedure is implemented sequentially via the following five steps, detailing the explicit matrix algebra that guarantees robustness.

\subsubsection*{Step 1: The First-Stage AR(1) (Purging the Predictor)}
We first isolate the innovations of the highly persistent predictor $X_t$ by estimating a first-order autoregression via OLS:
\begin{equation}
    X_t = \theta + \rho X_{t-1} + \varepsilon_t, \quad t=1,\dots,T
\end{equation}
The OLS estimators for the intercept and persistence parameter are given by:
\begin{equation}
    \begin{pmatrix} \widehat{\theta} \\ \widehat{\rho} \end{pmatrix} = \left( \sum_{t=1}^T \begin{pmatrix} 1 \\ X_{t-1} \end{pmatrix} \begin{pmatrix} 1 & X_{t-1} \end{pmatrix} \right)^{-1} \sum_{t=1}^T \begin{pmatrix} 1 \\ X_{t-1} \end{pmatrix} X_t
\end{equation}
We then compute and save the estimated residuals, which act as our proxy for the pure financial shock:
\begin{equation}
    \widehat{\varepsilon}_t = X_t - \widehat{\theta} - \widehat{\rho} X_{t-1}
\end{equation}

\subsubsection*{Step 2: The Joint Sieve-AR(1) Estimation (Purging the Weather Curve)}
To estimate the nonparametric weather curve $m(W_{t-1})$ without absorbing any endogenous financial shocks, we set up a joint OLS regression. The return shock is decomposed as $U_t = \phi \varepsilon_t + z_t$. Substituting this into the return equation yields the full regression model:
\begin{equation}
    Y_t = \bm{P}_K(W_{t-1})' \bm{c} + \beta X_{t-1} + \phi \widehat{\varepsilon}_t + \text{error}_t
\end{equation}
Let $\bm{\Lambda}_t = \left( \bm{P}_K(W_{t-1})', X_{t-1}, \widehat{\varepsilon}_t \right)'$ denote the $(K+2) \times 1$ concatenated regressor vector. The joint OLS estimator $\widehat{\bm{\delta}} = \left( \widehat{\bm{c}}', \widehat{\beta}, \widehat{\phi} \right)'$ is defined as:
\begin{equation}
    \widehat{\bm{\delta}} = \left( \sum_{t=1}^T \bm{\Lambda}_t \bm{\Lambda}_t' \right)^{-1} \sum_{t=1}^T \bm{\Lambda}_t Y_t
\end{equation}
We extract the first $K$ elements of $\widehat{\bm{\delta}}$ to serve as our robust estimator for the Sieve coefficients, $\widehat{\bm{c}}$. 
\begin{itemize}
    \item \textbf{No Omitted Variable Bias:} By explicitly including $X_{t-1}$ in $\bm{\Lambda}_t$, the projection of $Y_t$ onto the Sieve basis is orthogonalized against the predictor, preventing $\widehat{\bm{c}}$ from spuriously absorbing the predictive signal $\beta_0 X_{t-1}$.
    \item \textbf{No Stambaugh Bias:} By including $\widehat{\varepsilon}_t$, the endogeneity correlation $\mathbb{E}[U_t \varepsilon_t] \neq 0$ is explicitly partialled out. The remaining error is strictly orthogonal to the predictor's innovations, thereby unbiasing the estimation.
\end{itemize}

\subsubsection*{Step 3: Redefining the Adjusted Return}
With a clean estimator $\widehat{\bm{c}}$ in hand, we isolate the financial components of the return by subtracting the estimated nonparametric weather impact. The adjusted dependent variable is defined as:
\begin{equation}
    \widetilde{Y}_t = Y_t - \bm{P}_K(W_{t-1})'\widehat{\bm{c}}
\end{equation}
Under the null hypothesis $H_0: \beta = \beta_0$, substituting the true DGP $Y_t = m(W_{t-1}) + \beta_0 X_{t-1} + U_t$ and the Sieve decomposition $m(W_{t-1}) = \bm{P}_K(W_{t-1})'\bm{c} + r_t$ (where $r_t$ is the approximation error) yields:
\begin{align}
    \widetilde{Y}_t - \beta_0 X_{t-1} &= m(W_{t-1}) + U_t - \bm{P}_K(W_{t-1})'\widehat{\bm{c}} \\
    &= \bm{P}_K(W_{t-1})'\bm{c} + r_t + U_t - \bm{P}_K(W_{t-1})'\widehat{\bm{c}} \\
    &= U_t + r_t - \bm{P}_K(W_{t-1})'(\widehat{\bm{c}} - \bm{c})
\end{align}

\subsubsection*{Step 4: The Sieve Weight Projection}
In standard Empirical Likelihood, the estimation error $(\widehat{\bm{c}} - \bm{c})$ would severely distort the asymptotic distribution. To neutralize this, we construct an orthogonalized weight function. Let $w(X_{t-1})$ be the raw weight (e.g., $\tanh(X_{t-1}/c_w)$). We project $w(X_{t-1})$ onto the Sieve basis:
\begin{equation}
    \widehat{\bm{\gamma}} = \left( \sum_{t=1}^T \bm{P}_K(W_{t-1}) \bm{P}_K(W_{t-1})' \right)^{-1} \sum_{t=1}^T \bm{P}_K(W_{t-1}) w(X_{t-1})
\end{equation}
The orthogonalized, centered weight $w^c_t$ is defined as the residual of this regression:
\begin{equation}
    w^c_t = w(X_{t-1}) - \bm{P}_K(W_{t-1})' \widehat{\bm{\gamma}}
\end{equation}
By the fundamental first-order conditions of OLS, the residuals are exactly orthogonal to the regressors in-sample:
\begin{equation}
    \sum_{t=1}^T \bm{P}_K(W_{t-1}) w^c_t \equiv \bm{0}
\end{equation}

\subsubsection*{Step 5: Algebraic Elimination in the EL Moment Condition}
For testing $H_0: \beta = \beta_0$, the EL moment condition is evaluated as:
\begin{equation}
    Z_t(\beta_0) = \left( \widetilde{Y}_t - \beta_0 X_{t-1} \right) w^c_t
\end{equation}
Substituting the expansion from Step 3:
\begin{equation}
    Z_t(\beta_0) = \Big( U_t + r_t - \bm{P}_K(W_{t-1})'(\widehat{\bm{c}} - \bm{c}) \Big) w^c_t
\end{equation}
Summing the moment conditions across the sample to evaluate the numerator of the EL statistic:
\begin{equation}
    \sum_{t=1}^T Z_t(\beta_0) = \sum_{t=1}^T (U_t + r_t) w^c_t - (\widehat{\bm{c}} - \bm{c})' \underbrace{\sum_{t=1}^T \bm{P}_K(W_{t-1}) w^c_t}_{=\; \bm{0} \text{ (from Step 4)}}
\end{equation}
The multi-dimensional Sieve estimation error $\widehat{\bm{c}} - \bm{c}$ is thereby perfectly annihilated algebraically. We are left with:
\begin{equation}
    \frac{1}{\sqrt{T}} \sum_{t=1}^T Z_t(\beta_0) = \frac{1}{\sqrt{T}} \sum_{t=1}^T U_t w^c_t + \underbrace{\frac{1}{\sqrt{T}} \sum_{t=1}^T r_t w^c_t}_{=\; o_p(1)}
\end{equation}
Because the condition is structurally stripped of estimation variance, the sample variance of $Z_t(\beta_0)$ accurately reflects the pure innovation variance. As a result, the Empirical Likelihood ratio statistic (where $\widehat{\lambda}$ is the Lagrange multiplier) achieves the standard chi-squared limit:
\begin{equation}
    \ell_T(\beta_0) = 2 \sum_{t=1}^T \log(1 + \widehat{\lambda} Z_t(\beta_0)) \xrightarrow{d} \chi^2_1
\end{equation}

\subsection{Relation to Fixed-Intercept Predictive EL}

The framework differs from fixed-intercept predictive EL in three fundamental aspects:
\begin{itemize}
    \item \textbf{Nuisance Dimensionality}: A fixed scalar intercept $\alpha$ is replaced by an infinite-dimensional nuisance function $m(W_{t-1})$, which requires a sieve approximation of dimension $K \to \infty$.
    \item \textbf{Orthogonalization Space}: Simple weight demeaning against the constant subspace $\text{span}\{\bm{1}\}$ is replaced by projection onto the $K$-dimensional sieve space $\text{span}\{\bm{P}_K(W_{t-1})\}$.
    \item \textbf{Purging and Exogeneity}: By jointly estimating the sieve basis alongside $X_{t-1}$ and $\widehat{\varepsilon}_t$, the nonparametric component is protected against omitted predictor signal and predictor-innovation endogeneity.
\end{itemize}

\section{The Growth Rate of the Sieve Dimension $K$}

To establish the asymptotic validity of our test statistic, we must place bounds on the growth rate of the tuning parameter $K$ relative to the sample size $T$.

\subsection{Lower Bound (Bias Elimination)}
Assume the true weather curve $m(W)$ has smoothness $s$ and that $W$ is $d$-dimensional. Sieve theory implies that the approximation error $r_t$ satisfies $\sup |r_t| = O_p(K^{-s/d})$. To ensure that the approximation bias does not distort the asymptotic distribution, we require $\sqrt{T} \sum_{t=1}^T r_t w^c_t = o_p(1)$, which translates to:
\begin{equation}
    \sqrt{T} K^{-s/d} \to 0 \quad \implies \quad K \gg T^{\frac{d}{2s}}
\end{equation}

\subsection{Upper Bound (Variance Control and Leakage Mitigation)}
Estimating $K$ parameters in OLS introduces in-sample variance inflation and potential "future leakage" bias. Although the $\alpha$-mixing property of the weather process $W_t$ with exponential decay prevents leakage from accumulating, the pure variance penalty scales as $K / \sqrt{T}$. To prevent variance inflation from inflating the denominator of the EL ratio, we require:
\begin{equation}
    \frac{K}{\sqrt{T}} \to 0 \quad \implies \quad K \ll T^{\frac{1}{2}}
\end{equation}

\subsection{Growth rate of K}
Combining these bounds yields the formal asymptotic growth condition:
\begin{equation}
    T^{\frac{d}{2s}} \ll K \ll T^{\frac{1}{2}}
\end{equation}
This rate window is nonempty when $s>d$.
If $W$ is scalar and $m(W)$ is twice-differentiable ($d=1$, $s=2$), the condition becomes $T^{1/4} \ll K \ll T^{1/2}$. In practice, setting $K \approx T^{1/3}$ satisfies these bounds, ensuring that both approximation bias and variance inflation vanish asymptotically.

\section{The Formal Independence Assumption}

To ensure the algebraic orthogonalization rigorously eliminates the nonparametric estimation bias, we must formally state the economic relationship between the financial market predictor and the exogenous weather shock. We impose the following two independence conditions on $W_{t-1}$:

\paragraph*{Assumption 5.1 (Strict Exogeneity to the Market Shock)} \textit{The weather variable $W_{t-1}$ is strictly exogenous to the financial market innovation $U_t$. Mathematically:}
\begin{equation}
    \mathbb{E}[U_t \mid \mathcal{F}_{t-1}, W_{t-1}, W_t, W_{t+1}, \dots ] = 0.
\end{equation}

\noindent \textbf{Remark.} This condition ensures that using the entire dataset to estimate the weather curve does not inadvertently absorb future stock market shocks. It shuts down the "future leakage" trap. Because climatic processes are governed by atmospheric physics, weather shocks are completely oblivious to tomorrow's stock market crashes.

\paragraph*{Assumption 5.2 (Cross-Sectional Independence from the Predictor)} \textit{The exogenous weather process $\{W_t\}$ is strictly independent of the financial predictor innovation process $\{v_t\}$. Consequently, $W_{t-1}$ and $X_{t-1}$ are independent random variables, implying:}
\begin{equation}
    \mathbb{E}[w(X_{t-1}) \mid W_{t-1}] = \mathbb{E}[w(X_{t-1})].
\end{equation}

\noindent \textbf{Remark.} This assumption fundamentally preserves the predictor's signal. Because $W_{t-1}$ and $X_{t-1}$ are independent, projecting the EL weight $w(X_{t-1})$ onto the weather basis $\bm{P}_K(W_{t-1})$ produces an $R^2$ of exactly zero asymptotically. The sieve projection structurally deletes the nonparametric estimation bias without extracting a single drop of the persistent signal from $X_{t-1}$.

\section{Wilks' Theorem for Mean Predictability}
\label{sec:mean-wilks}

This section gives a formal Wilks theorem for the mean-predictability test.  It is
useful to write the nuisance step in null-imposed profile form.  Let
\[
    \bm Y=(Y_1,\ldots,Y_T)',\qquad
    \bm X_-=(X_0,\ldots,X_{T-1})',
\]
and let \(\bm P\) be the \(T\times K\) matrix whose \(t\)-th row is
\(\bm P_K(W_{t-1})'\).  Define
\[
    \bm\Pi_K=\bm P(\bm P'\bm P)^{-1}\bm P',
    \qquad
    \bm M_K=\bm I_T-\bm\Pi_K.
\]
For each candidate value \(\beta\), define the null-imposed profile sieve
coefficient and the corresponding residual by
\begin{equation}
    \widehat{\bm c}_K(\beta)
    =(\bm P'\bm P)^{-1}\bm P'(\bm Y-\beta\bm X_-),
    \qquad
    \widehat{\bm u}(\beta)
    =\bm M_K(\bm Y-\beta\bm X_-).
    \label{eq:profile-sieve}
\end{equation}
Let \(w_t=w(X_{t-1})\), \(\bm w=(w_1,\ldots,w_T)'\), and define the
sieve-orthogonalized weight
\begin{equation}
    \bm w^c=\bm M_K\bm w,
    \qquad
    w_t^c=(\bm w^c)_t.
    \label{eq:sieve-weight}
\end{equation}
The scalar estimating equation is
\begin{equation}
    \widehat Z_t(\beta)
    =\widehat u_t(\beta)w_t^c.
    \label{eq:profile-score}
\end{equation}
The empirical likelihood and its log-likelihood ratio are
\begin{align}
    L_T(\beta)
    &=\sup\left\{
       \prod_{t=1}^T(T\pi_t):
       \pi_t\geq0,\quad
       \sum_{t=1}^T\pi_t=1,\quad
       \sum_{t=1}^T\pi_t\widehat Z_t(\beta)=0
       \right\},                                                    \label{eq:profile-el}\\
    \ell_T(\beta)
    &=-2\log L_T(\beta)
      =2\sum_{t=1}^T\log\{1+\widehat\lambda(\beta)
          \widehat Z_t(\beta)\},                                    \label{eq:profile-elr}
\end{align}
where \(\widehat\lambda(\beta)\) solves
\begin{equation}
    \sum_{t=1}^T
       \frac{\widehat Z_t(\beta)}
            {1+\widehat\lambda(\beta)\widehat Z_t(\beta)}=0.
    \label{eq:lambda-equation}
\end{equation}

For a vector \(\bm a=(a_1,\ldots,a_T)'\), write
\(\|\bm a\|_T^2=T^{-1}\sum_{t=1}^T a_t^2\).  Also write
\[
    \overline w_T=T^{-1}\sum_{t=1}^T w_t,
    \qquad
    \widetilde w_t=w_t-\overline w_T.
\]

\begin{assumption}[Persistent-predictor score conditions]
\label{ass:score-core}
The predictor satisfies
\(X_t=\theta+\rho_T X_{t-1}+\varepsilon_t\), where \(\rho_T\) belongs to
one of the admissible persistence classes
including the stationary, mildly integrated, local-to-unity,
integrated, and mildly explosive cases.  With respect to the filtration
\[
    \mathcal H_t
    =\sigma\{(z_s,\varepsilon_s):s\leq t\}
       \vee\sigma\{W_s:s\in\mathbb Z\},
\]
\(E(U_t\mid\mathcal H_{t-1})=0\),
\(\sup_tE|U_t|^{2+\delta}<\infty\) for some \(\delta>0\), and the
conditional-variance and weight conditions needed for the oracle score limit
hold.  In particular, for
\[
    A_T=\frac1{\sqrt T}\sum_{t=1}^T U_t\widetilde w_t,
    \qquad
    V_T=\frac1T\sum_{t=1}^T U_t^2\widetilde w_t^2,
\]
there exists an a.s. positive, possibly random, variable \(\eta^2\) such that
\begin{equation}
    A_T\ \xrightarrow{\mathrm{st}}\ MN(0,\eta^2),
    \qquad
    V_T\ \xrightarrow{p}\ \eta^2,
    \qquad
    \max_{1\leq t\leq T}|U_t\widetilde w_t|=o_p(\sqrt T).
    \label{eq:score-core}
\end{equation}
The raw weight is uniformly bounded and satisfies the stated saturation
condition.
\end{assumption}

\begin{assumption}[Weather process and sieve projection]
\label{ass:sieve-projection}
The entire weather process is independent of
\(\{X_0,z_t,\varepsilon_t:t\geq1\}\), and the strict-exogeneity condition in
Assumption~\ref{ass:score-core} holds after enlarging the filtration by the
entire weather path.  The first component of \(\bm P_K(W)\) is one.  After
centering the remaining basis functions in sample, write the corresponding
\(T\times(K-1)\) matrix as \(\bm R\), so that
\(\bm 1_T'\bm R=\bm0\) and
\(\operatorname{span}(\bm P)=\operatorname{span}(\bm1_T,\bm R)\).
There are constants \(0<c<C<\infty\) such that, with probability tending to
one,
\[
    c\leq\lambda_{\min}(T^{-1}\bm R'\bm R)
      \leq\lambda_{\max}(T^{-1}\bm R'\bm R)\leq C.
\]
Let \(\zeta_K=\max_{t\leq T}\|\bm R_t\|\), where \(\bm R_t'\) is row \(t\)
of \(\bm R\).  The following increasing-dimensional score bounds hold:
\begin{align}
    \left\|T^{-1}\bm R'\widetilde{\bm w}\right\|
        &=O_p\!\left(\sqrt{K/T}\right),                                  \label{eq:rate-Rw}\\
    \left\|T^{-1/2}\bm R'\bm U\right\|&=O_p(\sqrt K),                    \label{eq:rate-RU}\\
    \left\|T^{-1}\bm P'\bm U\right\|
        &=O_p\!\left(\sqrt{K/T}\right),                                  \label{eq:rate-PU}
\end{align}
where \(\widetilde{\bm w}=(\widetilde w_1,\ldots,\widetilde w_T)'\) and
\(\bm U=(U_1,\ldots,U_T)'\).  These bounds are implied, for example, by
process independence, exponential mixing of \(W_t\), uniformly bounded basis
moments, and the martingale moment condition in
Assumption~\ref{ass:score-core}.
\end{assumption}

\begin{assumption}[Approximation and growth rates]
\label{ass:approximation}
For some \(\bm c_K\),
\[
    m(W_{t-1})=\bm P_K(W_{t-1})'\bm c_K+r_{K,t},
\]
where
\[
    \|\bm r_K\|_T=O(a_K),\qquad
    \|\bm r_K\|_\infty=O(a_K),\qquad
    \sqrt T\,a_K\longrightarrow0.
\]
The sieve projector is stable in sup norm on the approximation residual, so
that \(\|\bm M_K\bm r_K\|_\infty=O(a_K)\).  Moreover,
\begin{equation}
    K/\sqrt T\longrightarrow0,
    \qquad
    \zeta_K\sqrt{K/T}\longrightarrow0.
    \label{eq:K-rates}
\end{equation}
For normalized spline bases, for which typically \(\zeta_K\asymp\sqrt K\),
the second condition in \eqref{eq:K-rates} is implied by
\(K/\sqrt T\to0\).  If \(a_K\asymp K^{-s/d}\), a sufficient rate window is
\(T^{d/(2s)}\ll K\ll T^{1/2}\), which is nonempty when \(s>d\).
\end{assumption}

\begin{assumption}[Nondegeneracy and convex hull]
\label{ass:convex-hull}
There exists \(\underline\eta>0\) such that
\(P(\eta^2\geq\underline\eta^2)=1\).  In addition,
\[
    P\left\{
      \min_{1\leq t\leq T}\widehat Z_t(\beta_0)<0<
      \max_{1\leq t\leq T}\widehat Z_t(\beta_0)
    \right\}\longrightarrow1.
\]
\end{assumption}

\begin{lemma}[Negligibility of the growing sieve projection]
\label{lem:projection-negligible}
Under Assumptions~\ref{ass:sieve-projection}--\ref{ass:approximation},
\begin{align}
    \|\bm w^c-\widetilde{\bm w}\|_T
       &=O_p\!\left(\sqrt{K/T}\right),                                   \label{eq:weight-L2}\\
    \max_{t\leq T}|w_t^c-\widetilde w_t|
       &=O_p\!\left(\zeta_K\sqrt{K/T}\right)=o_p(1),                    \label{eq:weight-max}\\
    \frac1{\sqrt T}\sum_{t=1}^T
       U_t(w_t^c-\widetilde w_t)
       &=O_p(K/\sqrt T)=o_p(1),                                           \label{eq:weight-score}\\
    \|\bm\Pi_K\bm U\|_T
       &=O_p\!\left(\sqrt{K/T}\right),                                  \label{eq:U-proj-L2}\\
    \max_{t\leq T}|(\bm\Pi_K\bm U)_t|
       &=O_p\!\left(\zeta_K\sqrt{K/T}\right)=o_p(1).                   \label{eq:U-proj-max}
\end{align}
\end{lemma}

\begin{proof}
Because the sieve contains a constant and \(\bm1_T'\bm R=\bm0\), residualizing
\(\bm w\) on \(\bm P\) is equivalent to residualizing
\(\widetilde{\bm w}\) on \(\bm R\).  Hence
\[
    \bm w^c=\widetilde{\bm w}-\bm R\widehat{\bm\delta}_K,
    \qquad
    \widehat{\bm\delta}_K
       =(\bm R'\bm R)^{-1}\bm R'\widetilde{\bm w}.
\]
The Gram-matrix condition and \eqref{eq:rate-Rw} imply
\(\|\widehat{\bm\delta}_K\|=O_p(\sqrt{K/T})\).  Therefore,
\[
    \|\bm R\widehat{\bm\delta}_K\|_T^2
       =\widehat{\bm\delta}_K'
          (T^{-1}\bm R'\bm R)\widehat{\bm\delta}_K
       =O_p(K/T),
\]
which proves \eqref{eq:weight-L2}; and
\[
    \max_t|\bm R_t'\widehat{\bm\delta}_K|
       \leq\zeta_K\|\widehat{\bm\delta}_K\|
       =O_p\!\left(\zeta_K\sqrt{K/T}\right),
\]
which proves \eqref{eq:weight-max}.  Moreover,
\[
    \frac1{\sqrt T}\bm U'(\bm w^c-\widetilde{\bm w})
      =-\left(T^{-1/2}\bm R'\bm U\right)'
          \widehat{\bm\delta}_K
      =O_p(\sqrt K)O_p\!\left(\sqrt{K/T}\right)
      =O_p(K/\sqrt T),
\]
giving \eqref{eq:weight-score}.

Finally, let
\(\widehat{\bm a}_U=(\bm P'\bm P)^{-1}\bm P'\bm U\).  The Gram condition
and \eqref{eq:rate-PU} imply
\(\|\widehat{\bm a}_U\|=O_p(\sqrt{K/T})\).  The same quadratic-form and
maximum-row arguments give \eqref{eq:U-proj-L2} and
\eqref{eq:U-proj-max}.
\end{proof}

\begin{lemma}[Asymptotic equivalence of the feasible and oracle scores]
\label{lem:score-equivalence}
Under Assumptions~\ref{ass:score-core}--\ref{ass:approximation}, under
\(H_0:\beta=\beta_0\),
\begin{align}
    S_T
      :=\frac1{\sqrt T}\sum_{t=1}^T\widehat Z_t(\beta_0)
      &=\frac1{\sqrt T}\sum_{t=1}^T U_t\widetilde w_t+o_p(1),             \label{eq:numerator-equivalence}\\
    \widehat V_T
      :=\frac1T\sum_{t=1}^T\widehat Z_t(\beta_0)^2
      &=\frac1T\sum_{t=1}^T U_t^2\widetilde w_t^2+o_p(1),                \label{eq:variance-equivalence}\\
    \max_{t\leq T}|\widehat Z_t(\beta_0)|&=o_p(\sqrt T).                \label{eq:max-score}
\end{align}
\end{lemma}

\begin{proof}
Under the null,
\[
    \bm Y-\beta_0\bm X_-
       =\bm P\bm c_K+\bm r_K+\bm U,
    \qquad
    \widehat{\bm u}(\beta_0)
       =\bm M_K(\bm U+\bm r_K).
\]
Because \(\bm M_K\) is symmetric and idempotent and
\(\bm w^c=\bm M_K\bm w\),
\begin{align*}
    S_T
      &=T^{-1/2}\widehat{\bm u}(\beta_0)'\bm w^c\\
      &=T^{-1/2}(\bm U+\bm r_K)'\bm M_K^2\bm w\\
      &=T^{-1/2}(\bm U+\bm r_K)'\bm w^c.
\end{align*}
Lemma~\ref{lem:projection-negligible} gives
\(T^{-1/2}\bm U'\bm w^c
  =T^{-1/2}\bm U'\widetilde{\bm w}+o_p(1)\).  Also, by Cauchy--Schwarz and
the contraction property of an orthogonal projection,
\[
    \left|T^{-1/2}\bm r_K'\bm w^c\right|
       \leq\sqrt T\,\|\bm r_K\|_T\|\bm w^c\|_T
       \leq\sqrt T\,\|\bm r_K\|_T\|\bm w\|_T
       =o_p(1).
\]
This proves \eqref{eq:numerator-equivalence}.

For the second-moment result, define
\[
    e_{u,t}=\widehat u_t(\beta_0)-U_t,
    \qquad
    e_{w,t}=w_t^c-\widetilde w_t.
\]
Then
\[
    \bm e_u=-\bm\Pi_K\bm U+\bm M_K\bm r_K,
\]
so that, by Lemma~\ref{lem:projection-negligible} and projection contraction,
\[
    \|\bm e_u\|_T
       =O_p\!\left(\sqrt{K/T}+a_K\right)=o_p(1).
\]
Moreover, \(\max_t|e_{w,t}|=o_p(1)\),
\(\max_t|w_t^c|=O_p(1)\), and \(T^{-1}\sum_tU_t^2=O_p(1)\).  Since
\[
    \widehat Z_t(\beta_0)-U_t\widetilde w_t
       =e_{u,t}w_t^c+U_te_{w,t},
\]
it follows that
\begin{align*}
    \frac1T\sum_{t=1}^T
       \{\widehat Z_t(\beta_0)-U_t\widetilde w_t\}^2
    &\leq
       2\max_t|w_t^c|^2\|\bm e_u\|_T^2
       +2\max_t|e_{w,t}|^2\frac1T\sum_{t=1}^TU_t^2\\
    &=o_p(1).
\end{align*}
Equation \eqref{eq:variance-equivalence} now follows from Cauchy--Schwarz.

Finally, \(\max_t|U_t|=o_p(\sqrt T)\) follows from the uniform
\((2+\delta)\)-moment bound.  Together with
\eqref{eq:U-proj-max}, the sup-norm stability of \(\bm M_K\bm r_K\), and
\(\max_t|w_t^c|=O_p(1)\), this proves \eqref{eq:max-score}.
\end{proof}

\begin{theorem}[Wilks phenomenon for sieve-orthogonalized mean EL]
\label{thm:sieve-wilks}
Suppose Assumptions~\ref{ass:score-core}--\ref{ass:convex-hull} hold.  Then,
under \(H_0:\beta=\beta_0\),
\begin{equation}
    \ell_T(\beta_0)\ \xrightarrow{d}\ \chi_1^2.
    \label{eq:wilks-result}
\end{equation}
The conclusion holds for each predictor-persistence class covered by the high-level
score condition.
\end{theorem}

\begin{proof}
By Lemma~\ref{lem:score-equivalence} and
Assumption~\ref{ass:score-core},
\[
    S_T\ \xrightarrow{\mathrm{st}}\ MN(0,\eta^2),
    \qquad
    \widehat V_T\ \xrightarrow{p}\eta^2.
\]
Consequently,
\begin{equation}
    \frac{S_T}{\sqrt{\widehat V_T}}
       \ \xrightarrow{\mathrm{st}}\ N(0,1),
    \qquad
    \frac{S_T^2}{\widehat V_T}
       \ \xrightarrow{d}\chi_1^2.
    \label{eq:self-normalized-limit}
\end{equation}

It remains to connect the self-normalized statistic to empirical likelihood.
Write \(\widehat Z_t=\widehat Z_t(\beta_0)\) and
\(\widehat\lambda=\widehat\lambda(\beta_0)\).  The convex-hull condition,
\(S_T=O_p(1)\), \(\widehat V_T\to_p\eta^2>0\), and
\(\max_t|\widehat Z_t|=o_p(\sqrt T)\) imply the standard scalar-EL bounds
\[
    \widehat\lambda=O_p(T^{-1/2}),
    \qquad
    \max_{t\leq T}|\widehat\lambda\widehat Z_t|=o_p(1).
\]
Expanding the multiplier equation \eqref{eq:lambda-equation} gives
\[
    0=\sum_{t=1}^T\widehat Z_t
      -\widehat\lambda\sum_{t=1}^T\widehat Z_t^2
      +o_p(\sqrt T),
\]
and hence
\begin{equation}
    \widehat\lambda
       =\frac{\sum_{t=1}^T\widehat Z_t}
              {\sum_{t=1}^T\widehat Z_t^2}
        +o_p(T^{-1/2}).
    \label{eq:lambda-expansion}
\end{equation}
A second-order expansion of \eqref{eq:profile-elr}, together with
\(\max_t|\widehat\lambda\widehat Z_t|=o_p(1)\), yields
\begin{align*}
    \ell_T(\beta_0)
      &=2\widehat\lambda\sum_{t=1}^T\widehat Z_t
        -\widehat\lambda^2\sum_{t=1}^T\widehat Z_t^2+o_p(1)\\
      &=\frac{\left(T^{-1/2}\sum_{t=1}^T\widehat Z_t\right)^2}
              {T^{-1}\sum_{t=1}^T\widehat Z_t^2}
        +o_p(1)\\
      &=\frac{S_T^2}{\widehat V_T}+o_p(1).
\end{align*}
The result follows from \eqref{eq:self-normalized-limit} and Slutsky's
theorem.
\end{proof}

\begin{corollary}[Test of no mean predictability]
\label{cor:mean-test}
For \(H_0:\beta=0\), reject at asymptotic level \(\alpha\) whenever
\[
    \ell_T(0)>\chi^2_{1,1-\alpha}.
\]
Equivalently, the asymptotic \(p\)-value is
\(1-F_{\chi_1^2}\{\ell_T(0)\}\).  A confidence set for \(\beta\) is obtained
by inversion:
\[
    \mathcal C_{1-\alpha}
      =\{\beta:\ell_T(\beta)\leq\chi^2_{1,1-\alpha}\}.
\]
\end{corollary}

\begin{remark}[Relation to the unrestricted two-stage estimator]
\label{rem:unrestricted-chat}
The profile construction \eqref{eq:profile-sieve} is the cleanest route to
Theorem~\ref{thm:sieve-wilks}: at \(\beta_0\), the fitted nuisance error is
exactly \(\bm\Pi_K(\bm U+\bm r_K)\), whose prediction norm is
\(O_p(\sqrt{K/T}+a_K)\).  If one instead retains a separately estimated
unrestricted coefficient \(\widehat{\bm c}\), the same theorem continues to
hold provided one proves the additional generated-nuisance condition
\[
    \frac1T\sum_{t=1}^T
       \left\{\bm P_K(W_{t-1})'
              (\widehat{\bm c}-\bm c_K)\right\}^2(w_t^c)^2=o_p(1)
\]
and the corresponding maximum-term bound.  Those statements do not follow
from the algebraic identity \(\bm P'\bm w^c=\bm0\) alone.
\end{remark}

\begin{remark}[Why the same argument is not immediately a quantile theorem]
\label{rem:quantile}
For a quantile score \(\psi_\tau(u)=\tau-\mathbf1\{u<0\}\), the nuisance error
enters nonlinearly:
\[
  \psi_\tau\!\left(U_{t,\tau}+r_{K,t}
       -\bm P_K(W_{t-1})'(\widehat{\bm c}_\tau-\bm c_{K,\tau})\right).
\]
Therefore \(\bm P'\bm w^c=\bm0\) does not algebraically cancel the first-stage
error.  A first-order expansion instead produces a density-weighted term of
the form
\[
  \sum_{t=1}^T f_{t,\tau}(0)
       \bm P_K(W_{t-1})w_t^c,
\]
which is not generally zero.  A quantile extension would require either a
restrictive condition making \(f_{t,\tau}(0)\) constant with respect to the
predictor, or a new density-weighted orthogonalization and additional
nonparametric-rate arguments.  Thus Theorem~\ref{thm:sieve-wilks} is a
complete and natural endpoint for the mean paper, whereas a general quantile
extension is a separate methodological problem.
\end{remark}

\section{Conclusion}

When researchers use standard predictive regressions, they are often forced to assume a constant intercept to maintain mathematical tractability. Consequently, other exogenous macro-environmental factors are relegated to the error term. While this preserves the null hypothesis, it can inflate residual variance and weaken the test's power to detect financial predictability. The Orthogonalized Sieve-EL estimator removes nonlinear environmental noise through sieve profiling and protects the persistent-predictor signal through weight orthogonalization, while retaining standard empirical-likelihood calibration under the stated score conditions.

\end{document}
